
\documentclass{article}
\usepackage{colortbl}
\usepackage{makecell}
\usepackage{multirow}
\usepackage{supertabular}

\begin{document}

\newcounter{utterance}

\centering \large Interaction Transcript for game `cladder', experiment `full\_v1.5\_default', episode 4867 with Qwen3.5{-}9B.
\vspace{24pt}

{ \footnotesize  \setcounter{utterance}{1}
\setlength{\tabcolsep}{0pt}
\begin{supertabular}{c@{$\;$}|p{.15\linewidth}@{}p{.15\linewidth}p{.15\linewidth}p{.15\linewidth}p{.15\linewidth}p{.15\linewidth}}
    \# & \multicolumn{2}{c}{Player} && \multicolumn{2}{c}{Game Master} \\
    \hline

    \theutterance \stepcounter{utterance}  
    & & & \multicolumn{4}{p{0.6\linewidth}}{
        \cellcolor[rgb]{0.9,0.9,0.9}{
            \makecell[{{p{\linewidth}}}]{
                \texttt{\tiny{[P1$\langle$GM]}}
                \texttt{You are an expert in causal inference. The following question is not a typical commonsense query, but rather a meticulously designed question created by a professor specializing in causal inference, intended to assess the students' mastery of the course content.} \\
\texttt{Imagine a self{-}contained, hypothetical world with only the following conditions, and without any unmentioned factors or causal relationships: Alice waking up has a direct effect on rainfall. Traffic has a direct effect on rainfall. The overall probability of waking up late is 9\%. The probability of waking up on time and high rainfall is 53\%. The probability of waking up late and high rainfall is 8\%. Is the chance of high rainfall smaller when observing waking up late?} \\
\texttt{Start your answer with 'Yes' or 'No', followed by additional reasoning or evidence to support your explanation.} \\
            }
        }
    }
     \\ \\

    \theutterance \stepcounter{utterance}  
    & \multicolumn{4}{p{0.6\linewidth}}{
        \cellcolor[rgb]{0.561,0.737,0.561}{
            \makecell[{{p{\linewidth}}}]{
                \texttt{\tiny{[P1$\rangle$GM]}}
                \texttt{No} \\
\\ 
\texttt{**Reasoning:**} \\
\\ 
\texttt{To determine whether the chance of high rainfall is smaller when observing waking up late, we must compare the marginal probability of high rainfall, \$P(\textbackslash{}text\{High Rain\})\$, with the conditional probability of high rainfall given waking up late, \$P(\textbackslash{}text\{High Rain\} \textbackslash{}mid \textbackslash{}text\{Late\})\$.} \\
\\ 
\texttt{1.  **Calculate the marginal probability of high rainfall, \$P(\textbackslash{}text\{High Rain\})\$:**} \\
\texttt{    The problem states the total probability of waking up late is \$P(\textbackslash{}text\{Late\}) = 9\textbackslash{}\%\$. Consequently, the probability of waking up on time is \$P(\textbackslash{}text\{On Time\}) = 100\textbackslash{}\% {-} 9\textbackslash{}\% = 91\textbackslash{}\%\$.} \\
\texttt{    We are given the joint probability of waking up on time and high rainfall: \$P(\textbackslash{}text\{On Time\} \textbackslash{}cap \textbackslash{}text\{High Rain\}) = 53\textbackslash{}\%\$.} \\
\texttt{    We are given the joint probability of waking up late and high rainfall: \$P(\textbackslash{}text\{Late\} \textbackslash{}cap \textbackslash{}text\{High Rain\}) = 8\textbackslash{}\%\$.} \\
\\ 
\texttt{    Using the law of total probability:} \\
\texttt{    \$\$P(\textbackslash{}text\{High Rain\}) = P(\textbackslash{}text\{On Time\} \textbackslash{}cap \textbackslash{}text\{High Rain\}) + P(\textbackslash{}text\{Late\} \textbackslash{}cap \textbackslash{}text\{High Rain\})\$\$} \\
\texttt{    \$\$P(\textbackslash{}text\{High Rain\}) = 53\textbackslash{}} \\
            }
        }
    }
    & & \\ \\

    \theutterance \stepcounter{utterance}  
    & & & \multicolumn{2}{p{0.3\linewidth}}{
        \cellcolor[rgb]{0.9,0.9,0.9}{
            \makecell[{{p{\linewidth}}}]{
                \texttt{\tiny{[GM$|$GM]}}
                \texttt{No} \\
            }
        }
    }
    & & \\ \\

    \theutterance \stepcounter{utterance}  
    & & & \multicolumn{2}{p{0.3\linewidth}}{
        \cellcolor[rgb]{0.9,0.9,0.9}{
            \makecell[{{p{\linewidth}}}]{
                \texttt{\tiny{[GM$|$GM]}}
                \texttt{no} \\
            }
        }
    }
    & & \\ \\

    \theutterance \stepcounter{utterance}  
    & & & \multicolumn{2}{p{0.3\linewidth}}{
        \cellcolor[rgb]{0.9,0.9,0.9}{
            \makecell[{{p{\linewidth}}}]{
                \texttt{\tiny{[GM$|$GM]}}
                \texttt{game\_result = WIN} \\
            }
        }
    }
    & & \\ \\

\end{supertabular}
}

\end{document}
